\chapter{Thu Nhập Tăng, Kỷ Luật Phải Tăng}
\markboth{Thu Nhập Tăng, Kỷ Luật Phải Tăng}{As Income Rises, Discipline Must Rise}

\section{Tiền Nhiều Hơn Nhưng Vẫn Trống Tay}
\begin{truyen}{Khi Thu Nhập Tăng Mà Cảm Giác Thiếu Không Giảm}{When Income Rises but Shortage Remains}
\chuhoa{S}{au một thời gian dạy thêm ổn định, người thầy thấy thu nhập của mình khá hơn trước. Anh tưởng rằng từ đây áp lực tài chính sẽ nhẹ đi.}
\textit{(After a period of stable tutoring, the teacher saw his income improve. He thought financial pressure would now begin to ease.)}

Nhưng điều lạ là cuối tháng vẫn thấy thiếu. Không còn thiếu như cũ, nhưng vẫn có cảm giác tiền đi rất nhanh, còn sự yên tâm thì chưa đến. Anh ngồi nhìn lại mới hiểu: thu nhập tăng lên đã kéo theo một loạt mong muốn vốn trước đây còn bị nén xuống. Những thứ từng bị xem là không cần bỗng đổi tên thành phần thưởng, thành tiện nghi, thành điều xứng đáng.
\textit{(Yet strangely, he still felt short at the end of the month. Not short in the same way as before, but money still seemed to leave quickly while peace of mind had not yet arrived. Looking back, he understood: higher income had pulled upward a whole set of desires that had once been kept down. Things previously considered unnecessary were suddenly renamed rewards, comforts, and well-deserved purchases.)}

Anh phát hiện ra một sự thật không dễ nghe: nhiều người không nghèo chỉ vì kiếm ít, mà còn vì mỗi lần kiếm khá hơn thì nhu cầu cũng phình ra theo. Nếu không tăng kỷ luật cùng với thu nhập, người ta rất dễ chạy mãi mà vẫn không tới cảm giác đủ.
\textit{(He discovered a difficult truth: many people are not poor only because they earn little, but because each improvement in income makes their wants expand as well. If discipline does not rise together with income, one may keep running and still never reach the feeling of enough.)}
\end{truyen}

\begin{baihoc}
Thu nhập tăng là cơ hội, không phải giấy phép để thả lỏng mình.
\textit{(Higher income is an opportunity, not permission to loosen self-discipline.)}
\end{baihoc}

\begin{ghinhoanh}
More income is not automatic progress.

Discipline decides whether money stays.
\end{ghinhoanh}

\begin{tuhocnhanh}
1. Thu nhập tăng gần đây đã kéo theo những mong muốn nào của em? \textit{Which desires have risen with your recent increase in income?}

2. Em có đang tưởng thu nhập khá hơn đồng nghĩa với việc được dễ dãi hơn không? \textit{Are you mistaking better income for permission to be more careless?}
\end{tuhocnhanh}
\ngancach

\section{Gia Thư Tăng Quốc Phiên}
\begin{truyen}{Khi Gia Đình Suy Vì Quen Xa Hoa}{When Families Decline Through Luxury}
\chuhoa{T}{ăng Quốc Phiên nhiều lần răn con cháu phải giữ nếp thanh đạm. Ông hiểu một điều mà người đi lên thường rất dễ quên: cái nguy lớn không hẳn nằm ở lúc chưa có gì, mà nằm ở lúc vừa có một chút rồi bắt đầu xem sự xa hoa là bình thường.}
\textit{(Zeng Guofan repeatedly instructed his descendants to keep a simple life. He understood something that those rising in life often forget: the greatest danger does not always lie in having nothing, but in gaining a little and then treating luxury as normal.)}

Khi một gia đình quen tiêu theo mức cao hơn khả năng thực sự, chỉ cần một biến động nhỏ là mọi thứ xáo trộn. Cái khó không còn nằm ở vật chất trước mắt, mà ở việc lòng người đã quen rộng tay nên rất khó quay về nếp cũ.
\textit{(When a family becomes accustomed to spending above its true capacity, even a small shock can disrupt everything. The difficulty then lies not merely in material conditions, but in the fact that the heart has grown used to indulgence and finds it hard to return to older habits.)}

Lời dạy của người xưa vì thế rất hợp với người đang cố gắng đi lên: nghèo nhất thời chưa đáng sợ bằng mất nề nếp. Vì nề nếp mất rồi, tiền vào bao nhiêu cũng không thấy đủ.
\textit{(That ancient advice remains fitting for anyone trying to move upward: temporary poverty is less dangerous than the loss of discipline. Once discipline is lost, no amount of money feels sufficient.)}
\end{truyen}

\begin{baihoc}
Đi lên không khó bằng đi lên mà vẫn giữ được nếp sống lành mạnh.
\textit{(Rising in life is not as difficult as rising while still preserving a healthy way of living.)}
\end{baihoc}

\begin{ghinhoanh}
Wealth without discipline turns fragile.

Simple habits protect rising families.
\end{ghinhoanh}

\begin{tuhocnhanh}
1. Nề nếp tài chính nào của em đang có dấu hiệu lỏng đi? \textit{Which financial discipline in your life is starting to loosen?}

2. Điều gì em cần giữ dù thu nhập có tăng hơn nữa? \textit{What must you preserve even if your income continues to rise?}
\end{tuhocnhanh}
\ngancach

\section{Chia Tiền Trước Khi Tiền Biến Mất}
\begin{truyen}{Lương Cứng Một Ngăn, Tiền Dạy Thêm Một Ngăn}{Separate the Streams Before They Disappear}
\chuhoa{S}{au nhiều tháng loay hoay, người thầy bắt đầu đổi cách giữ tiền. Lương dạy học được dùng cho phần khung của gia đình. Tiền dạy thêm không còn được xem là khoản để tiêu thoải mái nữa.}
\textit{(After many months of struggling, the teacher changed the way he handled money. His teaching salary was reserved for the family’s basic structure. Tutoring income was no longer treated as money for free spending.)}

Anh chia tiền ngay khi vừa nhận được: một phần cho học phí con, một phần cho quỹ dự phòng, một phần để dành dài hơn, phần còn lại mới tính tới sinh hoạt linh hoạt. Cách chia ấy không làm anh giàu lên tức thì, nhưng nó làm tiền bớt chảy tràn như nước không có bờ.
\textit{(He divided the money as soon as it arrived: one part for the children’s school fees, one for emergencies, one for longer-term saving, and only the remainder for flexible daily use. That system did not make him rich immediately, but it prevented money from flowing away like water without banks.)}

Anh hiểu ra rằng thu nhập tăng chỉ phát huy giá trị khi nó được dẫn vào đúng chỗ. Nếu cứ để tất cả nằm chung trong một cảm giác ``mình đã vất vả nên được quyền tiêu'', thì cuối cùng mọi đồng tiền đều biến mất trong lý do nghe rất hợp tình.
\textit{(He came to understand that higher income becomes valuable only when directed into proper places. If everything remains mixed together under the feeling that ``I worked hard, so I deserve to spend,'' then every coin eventually disappears behind reasons that sound perfectly reasonable.)}
\end{truyen}

\begin{baihoc}
Tiền chỉ trở thành chỗ dựa khi được phân vai rõ ràng.
\textit{(Money becomes support only when its roles are clearly assigned.)}
\end{baihoc}

\begin{ghinhoanh}
Assign each coin a role.

Unplanned money disappears fastest.
\end{ghinhoanh}

\begin{tuhocnhanh}
1. Em có đang trộn tất cả nguồn thu vào một chỗ không? \textit{Are you mixing all income streams together?}

2. Nếu phải chia tiền ngay hôm nay, em sẽ chia theo bốn ngăn nào? \textit{If you had to divide money today, what four categories would you use?}
\end{tuhocnhanh}

\begin{turanminh}
1. Tôi phải nhớ rằng lương và tiền dạy thêm không sinh ra để chạy theo tiện nghi mới, mà để dựng lại tương lai chắc hơn cho gia đình.

2. Thu nhập tăng mà lòng ham cũng tăng theo thì mình chỉ đổi một kiểu thiếu này sang một kiểu thiếu khác.

3. Tôi phải chia tiền ngay khi nhận được, nếu không những lý do ``vất vả nên được quyền tiêu'' sẽ cuốn sạch mọi dự định tốt.
\end{turanminh}

\begin{dayhoc}
1. Học sinh cần hiểu sớm rằng kiếm được nhiều hơn chưa chắc sống vững hơn nếu không có kỷ luật.

2. Có thể dùng chương này để răn các em về thói quen tiêu hết tiền tiêu vặt hay tiêu hết phần thưởng ngay khi vừa có.

3. Một câu chốt cho lớp là: ``Tiền tăng mà kỷ luật không tăng thì sự yên tâm vẫn không tới.''
\end{dayhoc}